\section{Fault Injection Implementation}
\label{sec:fault-injection-impl}

Consistent with the Chaos Engineering methodology defined in Chapter \ref{chap:cloud-computing}, we utilize Chaos Mesh to programmatically inject faults into our target microservices cluster. The experiments defined in Chapter \ref{chap:methodology} are implemented as Kubernetes Custom Resource Definitions (CRDs). To evaluate the diagnostic capabilities of our observability architecture against a wide spectrum of system failures, we implement three declarative fault injection profiles:

\subsection{Pod-Level Disruption (PodChaos)}
The first campaign simulates a hard hardware or runtime crash of a specific service container, forcing Kubernetes to perform self-healing and redeploy the pod. This profile tests the preservation of forensic container logs in Loki after the container is destroyed. The declarative ``Pod Failure'' scenario is codified in Listing \ref{code:pod-chaos}.

\begin{lstlisting}[language=Yaml, caption={Pod Chaos Manifest}, label={code:pod-chaos}]
apiVersion: chaos-mesh.org/v1alpha1
kind: PodChaos
metadata:
  name: frontend-failure
spec:
  action: pod-kill
  mode: one
  selector:
    namespaces:
      - default
    labelSelectors:
      "app": "frontend"
  duration: "60s"
\end{lstlisting}

\subsection{Network-Level Performance Degradation (NetworkChaos)}
The second campaign introduces high-cardinality performance anomalies, injecting synthetic delay on service-to-service communication paths. This tests the semantic linkage of distributed traces in Tempo and metric-level anomalies under cascading latency. The ``Network Latency'' manifest targeting the payment service is codified in Listing \ref{code:network-chaos}.

\begin{lstlisting}[language=Yaml, caption={Network Latency Chaos Manifest}, label={code:network-chaos}]
apiVersion: chaos-mesh.org/v1alpha1
kind: NetworkChaos
metadata:
  name: payment-delay
spec:
  action: delay
  mode: one
  selector:
    namespaces:
      - default
    labelSelectors:
      "app": "payment"
  delay:
    latency: "250ms"
    jitter: "50ms"
    correlation: "25"
  direction: to
  duration: "120s"
\end{lstlisting}

\subsection{Resource Saturation (StressChaos)}
The final campaign simulates resource starvation and operational overload, consuming computational CPU or virtual memory within container limits to assess metric alerts and Prometheus threshold violations. The ``CPU Stress'' manifest designed to saturate the shopping cart microservice is codified in Listing \ref{code:stress-chaos}.

\begin{lstlisting}[language=Yaml, caption={CPU Stress Chaos Manifest}, label={code:stress-chaos}]
apiVersion: chaos-mesh.org/v1alpha1
kind: StressChaos
metadata:
  name: cart-cpu-stress
spec:
  mode: one
  selector:
    namespaces:
      - default
    labelSelectors:
      "app": "cart"
  stressors:
    cpu:
      workers: 2
      load: 100
  duration: "180s"
\end{lstlisting}

This declarative multi-profile approach ensures that the ``stress'' state is deterministic, reproducible, and mathematically identical for both the Vanilla and PANOPTES scenarios.